\chapter{The source side: lean4lean grounding, source evaluation, the supported fragment}
\label{chap:source}

A correctness statement for erasure needs three things fixed before the erasure relation
itself can be stated: what the source language means, which source programs the claim is
about, and how the source environment is mirrored on the target side. The Rocq development
of [S] gets the first for free --- PCUIC is a formalised type theory with its own reduction
relation --- but Lean has no formal operational semantics one can cite. This chapter
describes the substitute the development builds. It is grounded in \emph{lean4lean}, Mario
Carneiro's formalisation of Lean's kernel: a source term \code{Lean.Expr} acquires meaning
through lean4lean's translation relation \TrExprS{} into the abstract syntax of a kernel
model, whose judgements are typing and definitional equality. Nothing in this repository
re-proves any of that theory. What this chapter contributes is a big-step evaluation
relation on \code{Lean.Expr} whose every reduction step discharges an obligation back into
that model, a decidable predicate naming exactly which programs are covered, and the
environment-level companion of the erasure relation, assembled from what an erasure run
actually registered.

Four modules do the work. \code{LeanToLambdaBox/SourceEval.lean} defines \SEval{}, weak
call-by-value big-step evaluation of \code{Lean.Expr}, parameterised by a set of permitted
reductions, by a kernel environment, and --- the sharpest departure from the Rocq setting
--- by the \emph{compiler-body table}, since the Lean eraser reads the body the code
generator compiles, not the kernel's defining equation (design tag N8).
\code{LeanToLambdaBox/Supported.lean} names the fragment: a shape condition on a term and
on every definition body delta-reachable from it, with one named exclusion per audited
coverage hole, plus a total fuel-indexed decision procedure proved sound against it --- this
is how every green rung of Chapter~\ref{chap:capstone} discharges its coverage obligation.
\code{LeanToLambdaBox/SpecEnv.lean} builds the Lean analogue of MetaRocq's
\code{erases\_global} from an erasure run's own bookkeeping.
\code{LeanToLambdaBox/Upstream.lean} writes down, as one auditable structure, the four facts
about lean4lean's kernel theory that this development needs and cannot prove, because
editing the pinned fork is not this repository's work.

Two habits of the development are visible throughout and are worth stating once. First, the
value arms of \SEval{} are keyed on the \emph{source theory's} classification of a constant
(constructor, inductive type name, plain constant), not on a name test and not on the
compiler table, and the module spends four theorems proving that what the arms thereby
exclude is a fact rather than a convention. Second, coverage is a predicate and not a
footnote: the eighteen constructors of the error type each name a concrete defect ---
including two shapes on which the shipping eraser panics and emits a wrong program that
still passes \code{peregrine validate} --- and one of them exists purely to prevent
\emph{vacuous} coverage, since a recursor spine has no source evaluation at all.

\section{The lean4lean grounding}
\label{sec:source:lean4lean}

Everything in this chapter is stated relative to a kernel model imported from the pinned
lean4lean checkout. None of lean4lean's own theory is a node of this blueprint; it is the
imported background, and the single definition node below fixes the names the rest of the
document refers to.

\begin{definition}[The lean4lean grounding]
\label{def:lean4lean-grounding}
\lean{Lean4Lean.TrExprS, Lean4Lean.VExpr, Lean4Lean.VEnv, Lean4Lean.VLCtx,
Lean4Lean.VEnv.HasType, Lean4Lean.VEnv.IsDefEqU, Lean4Lean.TrProj}
\leanok

A source term \code{Lean.Expr} carries no meaning of its own. It acquires one through
lean4lean's translation relation $\TrExprS\ env\ Us\ \Delta\ e\ v$ --- "the source term $e$,
in level scope $Us$ and local context $\Delta$, translates to the kernel term $v$" --- where
$v$ ranges over the kernel model's term syntax (de Bruijn, no metadata, no literals),
$\Delta$ over local contexts pairing free-variable identities with kernel data, and $env$
over kernel environments. A kernel environment stores a constant table, a set of defining
equations and a pattern table; it stores \emph{no declaration list}, which is why every
"this name is a constructor / a type former / a plain constant" predicate used below
exhibits a well-formed declaration list producing an environment \emph{below} $env$, and why
those predicates are monotone in the environment. Typing is $env \vdash v : T$ and
definitional equality is the untyped-witness judgement $env \vdash v_1 \equiv v_2$; a
separate relation is the projection clause of the translation.
The Lean names are \code{Lean4Lean.TrExprS}, \code{Lean4Lean.VExpr}, \code{Lean4Lean.VEnv},
\code{Lean4Lean.VLCtx}, \code{Lean4Lean.VEnv.HasType}, \code{Lean4Lean.VEnv.IsDefEqU} and
\code{Lean4Lean.TrProj}. The pinned fork is not \code{sorry}-free: it contains sixteen
\code{sorry} sites, so every statement of this chapter whose proof routes through the
translation or the two judgements inherits them (Chapter~\ref{chap:trust}).
\end{definition}

\begin{remark}
Three of the capstone's hypothesis binders live entirely in this imported theory: the
subject's translation at the empty level scope and empty local context, the observable's
translation of the produced value, and that value's typing at an inductive spine. The first
is a checked term at every rung; the other two stand unproved at every rung and are
assumption nodes of Chapter~\ref{chap:trust}.
\end{remark}

\section{Source evaluation: flags, the step obligation, the relation}
\label{sec:source:seval}

\begin{definition}[The source application spine]
\label{def:mkApps}
\lean{LeanToLambdaBox.mkApps, LeanToLambdaBox.mkApps_const_inj}
\leanok
For a source head $f$ and arguments $a_1,\dots,a_n$,
$\mkApps f\ [a_1,\dots,a_n] = (\dots((f\,a_1)\,a_2)\dots a_n)$ is the left-nested application
spine. It is the shape on which every head-keyed rule of the source evaluation matches, and
it is injective at a constant head: if
$\mkApps (\mathtt{const}\ c\ us)\ args = \mkApps (\mathtt{const}\ c'\ us')\ args'$ then
$c = c'$, $us = us'$ and $args = args'$.
The Lean names are \code{mkApps} and \code{mkApps\_const\_inj}, both in the
\code{LeanToLambdaBox} namespace. This is the \emph{source} spine builder on
\code{Lean.Expr}; the target-side spine builder on \lbox{} terms
(Definition~\ref{def:LBTerm-mkApps}) is a different declaration in a different namespace,
and the capstone's observable mentions both.
\end{definition}

\begin{definition}[Permitted reductions]
\label{def:SEvalFlags}
\lean{LeanToLambdaBox.SEvalFlags, LeanToLambdaBox.fullFlags, LeanToLambdaBox.deltaOnly,
LeanToLambdaBox.SEvalFlags.le_refl, LeanToLambdaBox.SEvalFlags.le_trans,
LeanToLambdaBox.deltaOnly_le_fullFlags}
\leanok
Six booleans --- \code{beta}, \code{delta}, \code{zeta}, \code{iota}, \code{proj},
\code{lit} --- say which reductions a source evaluation may use. They are preordered by
"enables no more than", componentwise implication; the order is reflexive and transitive,
and the delta-only point is below the point enabling all six. That top point is the flag
point the capstone's observable conjunct runs at, and the one every rung of
Chapter~\ref{chap:capstone} uses. Making the reductions a parameter is what lets one
inductive relation cover every fragment of the schedule.
The Lean names are \code{SEvalFlags}, \code{fullFlags}, \code{deltaOnly} and the three order
facts. The order itself is an anonymous \code{LE} instance and so is not cited by name.
\end{definition}

\begin{definition}[The definitional equality a step owes]
\label{def:StepDefeq}
\lean{LeanToLambdaBox.StepDefeq, LeanToLambdaBox.StepDefeq.le}
\leanok
\uses{def:lean4lean-grounding}
$\mathsf{StepDefeq}\ env\ Us\ \Delta\ e_1\ e_2$ holds when \emph{both} sides translate into
the kernel model and their translations are definitionally equal there: there exist $v_1$
and $v_2$ with $\TrExprS\ env\ Us\ \Delta\ e_i\ v_i$ and
$env \vdash v_1 \equiv v_2$. It is the obligation attached to one source reduction
$e_1 \longrightarrow e_2$, and it is what ties an untyped-looking source step back to the
kernel theory; the delta, iota and projection arms of the evaluation relation each carry
one. It is monotone in the environment.
The Lean names are \code{StepDefeq} and \code{StepDefeq.le}. Both translations are bound
\emph{existentially} on purpose: a consumer stepping from $e_1$ needs the reduct's
translation in order to continue, and a form quantified over given translations cannot
supply one.
\end{definition}

\begin{lemma}[The existential form entails the universal one]
\label{lem:StepDefeq-uniq_form}
\lean{LeanToLambdaBox.StepDefeq.uniq_form}
\leanok
\uses{def:StepDefeq, def:lean4lean-grounding}
In a well-formed environment and local context,
$\mathsf{StepDefeq}\ env\ Us\ \Delta\ e_1\ e_2$ implies that \emph{every} pair of
translations of $e_1$ and $e_2$ is definitionally equal. So nothing is lost by binding the
translations existentially.
The Lean name is \code{StepDefeq.uniq\_form}; the Lean statement carries the two
well-formedness premises explicitly, on the environment and on the local context.
\end{lemma}
\begin{proof}
\leanok
\uses{def:StepDefeq, asm:lean4lean-trust}
Unpack the recorded pair $v_1, v_2$. lean4lean's uniqueness theorem for translations, read
at the reflexive local-context equality, identifies any other translation $w_1$ of $e_1$
with $v_1$ and $v_2$ with $w_2$ up to definitional equality; two applications of
transitivity of the definitional-equality judgement chain the three equalities into
$env \vdash w_1 \equiv w_2$.
\end{proof}

\begin{lemma}[\dots and the converse fails]
\label{lem:forall_form_not_stepDefeq}
\lean{LeanToLambdaBox.forall_form_not_stepDefeq, LeanToLambdaBox.trExprS_bvar_nil_elim}
\leanok
\uses{def:StepDefeq, def:lean4lean-grounding}
At $e_1 = \mathtt{bvar}\ 0$ in the empty local context the universally quantified side
condition holds vacuously --- nothing translates a loose de Bruijn index there --- while
$\mathsf{StepDefeq}$ does not. Hence the two forms are genuinely different and the
existential one is the right choice.
The Lean name is \code{forall\_form\_not\_stepDefeq}; its statement is a conjunction of the
vacuous universal form and the negation of the existential one, and the auxiliary
\code{trExprS\_bvar\_nil\_elim} records the underlying fact.
\end{lemma}
\begin{proof}
\leanok
\uses{def:StepDefeq}
The translation relation has no arm producing a translation of a loose \code{bvar} in the
empty context: case analysis on the \code{bvar} arm's lookup premise leaves a \code{none}.
The universal form is therefore vacuously true. The existential form would have to supply
exactly such a translation, so it fails.
\end{proof}

\begin{definition}[The eliminator's source-side segmentation]
\label{def:CasesOnShape}
\lean{LeanToLambdaBox.CasesOnShape, LeanToLambdaBox.MajorPremiseAt,
LeanToLambdaBox.CasesOnShape.mono}
\leanok
\uses{def:lean4lean-grounding, def:isCasesOnName}
$\mathrm{CasesOnShape}\ env\ c\ I\ dp\ nm$ says that $c$ is \emph{the} \code{casesOn} of the
inductive type $I$, at the segmentation $I$'s own declaration block fixes: $c$ passes the
\code{casesOn} name test, its name prefix is $I$, its declared type peels $dp$ binders and
then takes a major premise headed by $I$ at every level instantiation, and some well-formed
declaration list below $env$ declares a block containing $I$ with
$dp = \mathit{nparams} + 1 + (\mathit{arity} - \mathit{nparams})$ and
$nm = |\mathrm{ctors}(I)|$. The arithmetic is read off the \emph{inductive declaration}, not
off the environment's pattern table, so it does not depend on how that table was populated;
and since a \code{casesOn} constant is an ordinary constant in the model, the clause pins the
meaning of the name rather than filtering programs. The predicate survives environment
extension.
The Lean names are \code{CasesOnShape}, the auxiliary \code{MajorPremiseAt} and
\code{CasesOnShape.mono}; monotonicity holds because the declaring list already sits below
$env$ and a declared constant stays declared.
\end{definition}

\begin{remark}
The block that this predicate's existential exhibits is tied to the block the iota rule of
the evaluation relation reads its parameter count off only in company with the block
uniqueness carried by Definition~\ref{def:UpstreamAsks}; the module's own doc comment says
so. The converse, anti-monotone direction is strictly harder and has no consumer.
\end{remark}

\begin{definition}[Source evaluation]
\label{def:SEval}
\lean{LeanToLambdaBox.SEval}
\leanok
\uses{def:StepDefeq, def:CasesOnShape, def:mkApps, def:SEvalFlags, def:CtorOf, def:IndInfo,
def:IndArity, def:ConstOrigin, def:InformativeInd, def:lean4lean-grounding}
$\SEval\ env\ bo\ Us\ fl\ \Delta\ e\ v$ --- "in the kernel environment $env$, with
compiler-body table $bo$, level scope $Us$, permitted reductions $fl$ and local context
$\Delta$, the source term $e$ evaluates to the value $v$" --- is weak call-by-value big-step
evaluation of \code{Lean.Expr}, with eleven arms.
\begin{itemize}
\item \emph{Values that are binders or types.} A lambda is a value; a sort and a Pi type are
      values (weak evaluation does not enter a binder, and both shapes erase to $\Box$).
\item \emph{Constructor and type-former spines.} A constructor spine is a value once its
      arguments are \emph{and the spine is within the constructor's arity},
      $|args| \le np + nfs[k]$, read off the model's constructor and block data; an inductive
      type name applied to arguments is a value.
\item \emph{Beta and zeta} (guarded by the corresponding flags). The function evaluates to a
      lambda, the argument to a value, and the substituted body evaluates; likewise for
      \code{let}.
\item \emph{Delta on a compiler body} (guarded by the delta flag): $bo\ c$ is
      $\mathrm{some}\ b$, a guard says that $c$ is at no \code{CasesOnShape}, the body is
      instantiated at the call site's levels, all arguments evaluate pointwise, the step's
      $\mathsf{StepDefeq}$ is asked \emph{at the applied instance}, and the unfolded
      application evaluates to $v$. Without the eliminator guard a saturated eliminator spine
      would have a second derivation, unfolding the eliminator instead of reducing it; the
      eraser's own dispatch never visits a \code{casesOn} body.
\item \emph{Iota} (guarded by the iota flag): the subject is
      $\mkApps (\mathtt{const}\ con\ us)\ (\mathit{pre} \mathbin{+\!\!+}
      \mathit{disc} :: \mathit{minors} \mathbin{+\!\!+} \mathit{extra})$ with a
      \code{CasesOnShape} at that segmentation, a plain-constant reading of the head, a
      constructor reading of the selected constructor, the block's arity data and the
      inductive type's informativity; the discriminant evaluates to a constructor spine, the
      pre-arguments, the minors and the over-application all evaluate call by value, and the
      selected minor applied to the constructor's fields (the block's parameters dropped)
      together with the over-application evaluates to the result, with a
      $\mathsf{StepDefeq}$ recording the kernel's own reduction at that instance.
\item \emph{Projection} (guarded by the projection flag): the discriminee evaluates to a
      constructor spine of the structure and spine position $np + i$ evaluates, with the
      constructor reading, the block's single-constructor arity data and the bound
      $np + i < |cargs|$.
\item \emph{Literals} (guarded by the literal flag): a literal evaluates by unfolding to its
      constructor form, the step both the translation relation and the eraser take.
\end{itemize}
The Lean name is \code{SEval}. The value arms are [S, Fig.~12]'s \code{value\_head}
classification, arity bound included, keyed on the source theory's own reading of
\code{Expr.const} rather than on a name test or on the compiler table. The informativity
premise of the iota arm is restriction N18 --- the source relations model only the
eliminations the target performs --- and the projection arm carries \emph{no} relevance
premise, because the erasure relation's projection rule carries it.
\end{definition}

\begin{remark}
The argument-wise premises of the spine arms index with \code{getElem!}, so an out-of-range
read silently yields a default. The arms carry the length equations that make most reads
in-range, but the constructor arm's arity bound reads the field-count list at the
constructor index with no accompanying bound on that index, so an out-of-range index would
degenerate the bound.
\end{remark}

\begin{lemma}[Widening along the flags]
\label{lem:SEval-mono}
\lean{LeanToLambdaBox.SEval.mono}
\leanok
\uses{def:SEval, def:SEvalFlags}
If $fl \le fl'$ then every $\SEval$ derivation at $fl$ is one at $fl'$, with the \emph{same}
value.
The Lean name is \code{SEval.mono}.
\end{lemma}
\begin{proof}
\leanok
\uses{def:SEval}
Induction on the derivation. Every arm is rebuilt with the same premises, the only change
being that each flag guard is transported along the corresponding component of $fl \le fl'$.
The environment-facing side conditions are untouched.
\end{proof}

\begin{lemma}[Stability under an environment extension]
\label{lem:SEval-le}
\lean{LeanToLambdaBox.SEval.le}
\leanok
\uses{def:SEval, def:CasesOnShape}
If $env \le env'$ and no \code{CasesOnShape} of $env'$ is new --- \ie{} every
$\mathrm{CasesOnShape}\ env'\ c\ I\ dp\ nm$ is already one of $env$ --- then every $\SEval$
derivation in $env$ is one in $env'$. The extra hypothesis is not slack: the delta arm's
eliminator guard is a \emph{negative} fact about the environment, and an extension declaring
a block named after a constant's prefix makes that constant a \code{casesOn} head it was not.
The Lean name is \code{SEval.le}.
\end{lemma}
\begin{proof}
\leanok
\uses{def:SEval, def:StepDefeq, def:CasesOnShape}
Induction on the derivation. The value arms transport by monotonicity of the constructor and
block readings; the stepping arms transport their $\mathsf{StepDefeq}$ by its own
monotonicity lemma; the delta arm's negative guard is transported by contraposing the extra
hypothesis.
\end{proof}

\begin{remark}
Neither of the two preceding lemmas has a consumer anywhere in the tree, and neither has the
delta-only flag point: every actual consumer works at a generic flag parameter or directly
at the full flag point. The module docstring's "widening axis" story --- a simulation proved
at a narrow point and transported --- describes a schedule the tree no longer follows.
\end{remark}

\section{What is not a value}
\label{sec:source:values}

The four results of this section prove that what the value arms exclude is a fact about the
source theory rather than a convention of the definition.

\begin{lemma}[A body-less plain constant heads no value]
\label{lem:untabled_const_not_value}
\lean{LeanToLambdaBox.untabled_const_not_value}
\leanok
\uses{def:SEval, def:mkApps, def:CasesOnShape, def:CtorOf, def:IndInfo}
If $bo\ c = \mathrm{none}$ and $c$ is neither a constructor, nor an inductive type name, nor
at any \code{CasesOnShape}, then no spine $\mkApps (\mathtt{const}\ c\ us)\ args$ evaluates
at all, at any flags. This is PCUIC's treatment of an axiom, and it is what keeps an
axiom-freedom premise out of the simulation: a run that reaches such a head simply has no
source derivation.
The Lean name is \code{untabled\_const\_not\_value}; the three exclusions are stated as
universally quantified negations over the model's classifications.
\end{lemma}
\begin{proof}
\leanok
\uses{def:SEval, def:mkApps}
Strong induction on the spine length, then case analysis on the derivation. The two value
arms are excluded by the classification hypotheses; the lambda, sort, Pi, literal,
projection and zeta arms are excluded because their subjects cannot be a constant-headed
spine, which is where the spine inversion and the injectivity of the spine at a constant
head are spent; beta peels one argument and appeals to the induction hypothesis; delta needs
$bo\ c$ to be \code{some}; iota needs a \code{CasesOnShape}.
\end{proof}

\begin{lemma}[An over-applied constructor spine has no value]
\label{lem:overapplied_ctor_not_value}
\lean{LeanToLambdaBox.OverApplied, LeanToLambdaBox.overapplied_ctor_not_value}
\leanok
\uses{def:SEval, def:mkApps, def:CtorOf, def:IndInfo}
Say $n$ arguments \emph{over-apply} $c$ when no classification of $c$ in the environment
gives it an arity that long: for every constructor reading of $c$ in a block with parameter
count $np$ and field counts $nfs$, one has $np + nfs[k] < n$. Then a constructor spine with
that many arguments has no source value. This is exactly what the constructor value arm's
arity bound buys: without it the arm would call a value a spine on which the target is stuck.
The Lean names are \code{OverApplied} and \code{overapplied\_ctor\_not\_value}; the theorem
additionally assumes that $c$ is untabled and at no \code{CasesOnShape}.
\end{lemma}
\begin{proof}
\leanok
\uses{def:SEval, lem:CtorOf-not_indInfo}
Case analysis on the derivation at the given spine. The constructor value arm contradicts
over-application directly; the type-former value arm is excluded because a constructor name
is not an inductive type name (Lemma~\ref{lem:CtorOf-not_indInfo}); the remaining arms are
excluded exactly as in the previous lemma.
\end{proof}

\begin{remark}
Over-application quantifies over \emph{every} classification because the constructor and
block readings each read a declaration list below the environment and nothing in this
statement says two such lists agree. The uniqueness conjuncts of
Definition~\ref{def:UpstreamAsks} would collapse the quantifier to a single bound.
\end{remark}

\begin{lemma}[A constructor is not a type name]
\label{lem:CtorOf-not_indInfo}
\lean{LeanToLambdaBox.CtorOf.not_indInfo, LeanToLambdaBox.CtorOf.constant_ctorResult,
LeanToLambdaBox.IsArity.piBody_sort}
\leanok
\uses{def:CtorOf, def:IndInfo, def:IsArity}
A name that is a constructor of some inductive type in the model is not also an inductive
type name of that model.
The Lean name is \code{CtorOf.not\_indInfo}, with the two auxiliaries
\code{CtorOf.constant\_ctorResult} and \code{IsArity.piBody\_sort}.
\end{lemma}
\begin{proof}
\leanok
\uses{def:CtorOf, def:IndInfo}
A constructor's declared type is registered by lean4lean's inductive-declaration staging
with a result headed by its type former, while a type former's declared type is an arity,
whose Pi-body is a sort. A sort is not a constructor spine, so the two readings cannot
coincide.
\end{proof}

\begin{remark}
This lemma is one of the facts \code{doc/upstream-asks.md} files for consolidation upstream.
It also discharges, inside this repository, the first conjunct of the origin field of
Definition~\ref{def:UpstreamAsks}, which the field's own doc comment records as redundant
and retains so that the field grants everything the upstream ask was filed for.
\end{remark}

\begin{lemma}[A stuck head stays a head]
\label{lem:seval_const_head_value}
\lean{LeanToLambdaBox.seval_const_head_value, LeanToLambdaBox.SEval.forallE_inv}
\leanok
\uses{def:SEval, def:mkApps, def:CasesOnShape}
At a head where delta cannot move ($bo\ c = \mathrm{none}$) and iota cannot fire (no
\code{CasesOnShape}), the value of a spine on $c$ is a spine on $c$ --- in particular it is
never a lambda, so no such spine is the function of a beta redex. Similarly, a Pi type's only
source value is itself.
The Lean names are \code{seval\_const\_head\_value} and \code{SEval.forallE\_inv}.
\end{lemma}
\begin{proof}
\leanok
\uses{def:SEval, def:mkApps}
Strong induction on the spine length with the statement generalised over the subject, then
case analysis. The two value arms return a spine on the same head; beta peels an argument
and appeals to the induction hypothesis; the remaining arms are excluded by the two
hypotheses or by the shape of their subjects.
\end{proof}

\begin{remark}
The Pi-type inversion is what makes the Pi value arm non-optional: without it a delta redex
with a Pi-typed argument has no derivation at all.
\end{remark}

\section{The compiler-body table}
\label{sec:source:bodies}

\begin{definition}[The compiler-body table's typing hypothesis]
\label{def:CompilerBodies}
\lean{LeanToLambdaBox.CompilerBodies, LeanToLambdaBox.CompilerBodies.le}
\leanok
\uses{def:lean4lean-grounding}
$\mathrm{CompilerBodies}\ lenv\ env\ bo$ says every tabled compiler body is kernel-typeable
at its constant's declared type: for each $c$ with $bo\ c = \mathrm{some}\ b$, the constant
is declared in the elaboration environment $lenv$, and $b$ together with $c$'s declared type
both translate, with the translated body having the translated type at $c$'s own level scope
in the empty local context. It is monotone in the environment, and it is checkable per
program by running lean4lean's checker on the pair.
The Lean names are \code{CompilerBodies} and \code{CompilerBodies.le}.
\end{definition}

\begin{remark}
This predicate is a \emph{hypothesis} at every consumer: it is the capstone binder
\code{hcb} and a binder of every rung, and it is discharged at the first rung only. At the
other rungs it is blocked upstream, because a translation at a projection node routes through
lean4lean's projection relation, which is entirely unproven at the pin. Design tag N8: the
eraser's input is the compiler-facing body --- after \code{csimp} rewriting and
\code{@[extern]} or \code{@[implemented\_by]} substitution --- not the kernel-checked one.
See the assumption node for \code{hcb} in Chapter~\ref{chap:trust}.
\end{remark}

\section{Three environment fixtures}
\label{sec:source:witnesses}

Three fixture namespaces build concrete kernel environments by hand and exhibit derivations
in them, showing that the arms and their side conditions are jointly inhabitable rather than
vacuous. They account for the great majority of \code{SourceEval.lean}'s roughly 378
declarations: term abbreviations, per-stage constant lookups proved by \code{rfl},
per-subterm typing and translation derivations, the extension chains between staged
environments, and the field proofs of two inductive-declaration well-formedness instances.
Only the three results below state something a reader needs.

\begin{proposition}[A tabled body that is not the kernel body]
\label{prop:DeltaWitness-seval_sel_applied}
\lean{LeanToLambdaBox.DeltaWitness.seval_sel_applied,
LeanToLambdaBox.DeltaWitness.seval_add_deltaOnly,
LeanToLambdaBox.DeltaWitness.seval_add_applied,
LeanToLambdaBox.DeltaWitness.selBody_ne_addBody}
\leanok
\uses{def:SEval, def:mkApps, def:SEvalFlags, def:StepDefeq}
In a two-constant fixture environment, one constant's tabled compiler body \emph{is} its
kernel body while the other's is \emph{not}; nevertheless the applied instance of the second
unfolds its tabled body and beta-reduces to a sort, with the delta arm's
$\mathsf{StepDefeq}$ discharged --- because the obligation is asked at the \emph{applied}
instance, where both bodies compute to the same value.
The Lean name is \code{DeltaWitness.seval\_sel\_applied}, with the two companion derivations
and the inequality of the two bodies.
\end{proposition}
\begin{proof}
\leanok
\uses{def:SEval, def:StepDefeq}
Apply the delta arm with the tabled body; its eliminator guard is a name test in this
fixture; the two arguments evaluate by the sort and Pi value arms, which is where those arms
are spent, since the Pi-typed argument has no other derivation; the step's definitional
equality chains the fixture's recorded defining equation with the two bodies' applied
computations; the continuation is two beta steps to a sort.
\end{proof}

\begin{remark}
This is the single most important witness in the file: it shows that the delta arm really
does model compiler bodies that differ from kernel bodies, which is restriction N8 made
concrete. One caveat: the fixture's environment is never shown to be a well-formed kernel
environment (it sets the pattern table to \code{False} and the defining equations to a
two-element predicate), so the correct reading is "the arm's side conditions are jointly
satisfiable", not "the arm fires in a well-formed environment". The obligation itself needs
only the translation and the definitional-equality judgement.
\end{remark}

\begin{proposition}[The branch rule fires]
\label{prop:NatWitness-seval_iota_fires}
\lean{LeanToLambdaBox.NatWitness.seval_iota_fires,
LeanToLambdaBox.NatWitness.seval_iota_overapplied_fires,
LeanToLambdaBox.NatWitness.nat_casesOnShape, LeanToLambdaBox.NatWitness.natDecl_wf,
LeanToLambdaBox.NatWitness.nat_wf'_cas, LeanToLambdaBox.NatWitness.nat_stuck_not_value,
LeanToLambdaBox.NatWitness.natEnv, LeanToLambdaBox.NatWitness.nat_indInfo}
\leanok
\uses{def:SEval, def:CasesOnShape, def:SEvalFlags, def:mkApps, def:IndInfo}
A hand-built \code{Nat} block, declared through lean4lean's inductive-declaration staging and
proved well-formed, supports a \code{CasesOnShape} at the segmentation "one argument before
the major premise, one minor per constructor", and the iota arm fires on it both at the
saturated instance and at one argument past the emitted node --- the shape the eraser applies
\emph{outside} the emitted branch node. The same fixture's body-less constant heads no value.
The Lean name is \code{NatWitness.seval\_iota\_fires}, at the full flag point.
\end{proposition}
\begin{proof}
\leanok
\uses{def:SEval, def:StepDefeq, lem:untabled_const_not_value}
Instantiate the iota arm with the fixture's own classification facts --- its constructor
reading, its block arity, its informativity and the plain-constant reading of the eliminator
--- and its \code{CasesOnShape}. The pre-argument, the discriminant and the minors evaluate
by the value arms; the step's definitional equality is one of two hand-recorded conversion
equations; the continuation is a beta step.
\end{proof}

\begin{remark}
Two caveats. Because the registered pattern shape of the staging puts the major premise
\emph{last} while \code{casesOn} takes it before its minors, the fixture's environment is the
staged one \emph{extended with two hand-written defining equations}, and well-formedness is
proved only for the staged environment; the satisfiability claim is correspondingly weaker
than \code{doc/trust.md} reads. Second, this fixture is consumed outside the file: a rung of
Chapter~\ref{chap:capstone} inhabits its spike-facts bundle from six of the fixture's facts,
and \code{natSpecEnv} (Proposition~\ref{prop:natEnv_elimCovered}) builds the \lbox{} side of
the same fixture.
\end{remark}

\begin{proposition}[The projection rule fires]
\label{prop:ProjWitness-seval_proj_fires}
\lean{LeanToLambdaBox.ProjWitness.seval_proj_fires, LeanToLambdaBox.ProjWitness.p_projStep,
LeanToLambdaBox.ProjWitness.p_trProj, LeanToLambdaBox.ProjWitness.pDecl_wf,
LeanToLambdaBox.ProjWitness.p_wf'}
\leanok
\uses{def:SEval, def:SEvalFlags, def:lean4lean-grounding}
A hand-built two-field structure with one constructor, declared through the staging so that
its iota rule is registered the way a translated environment registers one, supports the
projection arm at the full flag point: the second field of a constructor spine is selected,
with the step's definitional equality read off a recorded defining equation and the
projection premise of the translation relation supplied at the projection's expansion.
The Lean name is \code{ProjWitness.seval\_proj\_fires}.
\end{proposition}
\begin{proof}
\leanok
\uses{def:SEval, def:StepDefeq}
The discriminee evaluates to the constructor spine by the constructor value arm; the
fixture's constructor and block-arity facts fix the head and the parameter count; the field
index is in range; the step's obligation is the recorded projection step, built from the
defining equation and the projection witness of the translation.
\end{proof}

\begin{remark}
The same caveat as the \code{Nat} fixture applies: the environment is the staged one extended
with a defining equation, and well-formedness is proved only before that extension. Worth
adding is the measurement \code{doc/trust.md} records: at a projection the exposure to
lean4lean is \emph{inhabitation}, not axioms --- a projection in a real program has a
translation derivation only through a lean4lean lemma that is \code{sorry} at the pin, so the
projection arm is non-vacuous only on hand-built witnesses.
\end{remark}

\section{The supported fragment}
\label{sec:source:fragment}

\begin{definition}[The coverage holes, named]
\label{def:SupportError}
\lean{LeanToLambdaBox.SupportError}
\leanok
Eighteen constructors, one per audited coverage hole, each naming a concrete defect:
\code{sparseCasesOn}, \code{sideConditionElim}, \code{etaContractedMinor}, \code{strLit},
\code{machineNat}, \code{quotPrim}, \code{ioLike}, \code{implementedBy}, \code{mvar},
\code{propElimIntoData}, \code{underAppliedCtor}, \code{underAppliedElim},
\code{recursorHead}, \code{mdataSpine}, \code{projField}, \code{kernameCollision},
\code{unknownConst}, \code{outOfFuel}. The checker returns the \emph{name} of the hole it
found, which is what generates \code{doc/coverage.md}.
The Lean name is \code{SupportError}.
\end{definition}

\begin{remark}
The exclusions fall into three kinds. \emph{Eraser bugs}: the sparse-\code{casesOn} shape, on
which the shipping eraser panics and emits a wrong program that still passes
\code{peregrine validate} (finding F-SPARSE), and string literals, on which the literal
visitor panics and returns $\Box$. \emph{Modelling gaps where the eraser eta-expands and the
erasure relation has no eta rule}: under-applied constructors and eliminators,
eta-contracted minors, metadata-headed spines. \emph{Vacuity guards}: the recursor head,
because a recursor spine has \emph{no source evaluation at all}, so a program reaching one
would be covered vacuously --- restriction N21. Two exclusions must be attributed correctly.
\code{implementedBy} is \emph{never produced} by the checker and its doc comment says so; the
shape is closed instead by the configuration hypothesis \code{ConfigPinned}
(Definition~\ref{def:ConfigPinned}), whose \code{csimp} and \code{extern} settings differ
from the default \code{\#erase} configuration. And \code{quotPrim} is how \code{Quot}
primitives in computationally relevant positions (finding F-QUOT) leave the fragment.
\end{remark}

\begin{remark}
The measured exclusions on the five benchmark programs are: a recursor head at
\code{Eq.rec} through \code{Bool.noConfusion} (finding F-EQREC) on Sieve, BinaryTrees,
Quicksort and Fannkuch; the well-founded \code{Nat.div.go} and \code{Nat.modCore.go} route
and a sparse \code{casesOn} additionally on Quicksort; an eta-contracted minor at
\code{Decidable.casesOn} additionally on Fannkuch. \textbf{Arith alone is in the fragment.}
\end{remark}

\begin{definition}[Manifest lambda-telescopes]
\label{def:IsLamTelescope}
\lean{LeanToLambdaBox.IsLamTelescope, LeanToLambdaBox.isLamTelescopeB,
LeanToLambdaBox.isLamTelescopeB_iff, LeanToLambdaBox.IsLamTelescope.instantiate1}
\leanok
$e$ is a manifest lambda-telescope of depth at least $n$ when it is syntactically $n$ nested
lambdas. The iota fragment needs it: the eraser's intro branch eta-expands a non-lambda minor
premise and the erasure relation has no eta rule. Lean's \code{match} compiler emits minors
as explicit \code{fun a b => \dots}, so real pattern-matching code is inside the fragment and
hand-written eta-contracted minors are not. A decision procedure decides it, and the
predicate survives opening a binder with a free variable.
The Lean names are \code{IsLamTelescope}, the Boolean \code{isLamTelescopeB}, their
correspondence and the instantiation lemma.
\end{definition}

\begin{definition}[Saturation of a constructor occurrence]
\label{def:CtorSaturated}
\lean{LeanToLambdaBox.CtorSaturated, LeanToLambdaBox.ctorSaturatedB,
LeanToLambdaBox.ctorSaturatedB_iff, LeanToLambdaBox.ctorOf?}
\leanok
\uses{def:Witness-SourceTable}
A tabled constructor occurs applied to at least its parameters and its fields; the condition
is vacuous at every other head. This is restriction N19's constructor half. Under-application
is excluded because the eraser eta-expands it and pushes the supplied prefix under the new
binders, where weak evaluation never reaches it.
The Lean names are \code{CtorSaturated}, its Boolean twin, their correspondence, and the
table lookup \code{ctorOf?} that recovers a tabled constructor from its name prefix.
\end{definition}

\begin{remark}
The condition was measured to fire zero times over the thirteen reified tables. Its
constructor half disappears once finding F-ETA2 is repaired, since applied-form \lbox{}
evaluates a partially applied constructor spine natively.
\end{remark}

\begin{definition}[The four head exclusions]
\label{def:PlainHead}
\lean{LeanToLambdaBox.PlainHead, LeanToLambdaBox.quotPrimNames, LeanToLambdaBox.ioLikeRoots,
LeanToLambdaBox.isIoLike, LeanToLambdaBox.isSparseCasesOn, LeanToLambdaBox.isMatcherName,
LeanToLambdaBox.isRecursorName}
\leanok
\uses{def:SupportError, def:Witness-SourceTable}
Every head inside the fragment is \emph{plain}: not a \code{Quot} primitive (\code{Quot},
\code{Quot.mk}, \code{Quot.lift}, \code{Quot.ind} --- \code{Quot.sound} is deliberately
absent, being \code{Prop}-typed hence erased), not \code{IO}-like (a list of roots including
\code{IO}, \code{ST}, \code{Task}, \code{String}, the floats and the machine integers), not a
sparse \code{casesOn} auxiliary, and not a surviving matcher. Separately, a head is rejected
when it is a recursor name whose prefix is a tabled inductive.
The Lean name of the four-field structure is \code{PlainHead}; the name tests it is built
from are \code{quotPrimNames}, \code{isIoLike}, \code{isSparseCasesOn} and
\code{isMatcherName}, and the separate recursor test is \code{isRecursorName}.
\end{definition}

\begin{remark}
Two defects worth recording. The sparse-\code{casesOn} diagnostic is also reported for a
\code{casesOn} whose name prefix is simply \emph{not tabled}, while the constructor's doc
comment defines the hole as a name-shape defect; since the error names generate
\code{doc/coverage.md}, a coverage row may mean "untabled inductive". And the recursor test
fires only when the recursor's prefix is a tabled inductive, so a program mentioning
\code{Foo.rec} while \code{Foo} is absent from the table's inductive column but present in
its constant column falls through to the constant branch --- exactly the vacuous coverage the
exclusion exists to prevent. Neither has a soundness consequence for the theorems below.
\end{remark}

\begin{definition}[A head the table knows]
\label{def:KnownHead}
\lean{LeanToLambdaBox.KnownHead}
\leanok
\uses{def:IndInfo, def:CtorOf, def:Witness-SourceTable}
Three arms: a tabled inductive type together with the model's block data, a tabled
constructor together with the model's constructor reading, and a tabled constant that the
model contains. Each arm carries the model's own reading of the name, which is what lets a
consumer eliminate the two columns its run has already excluded. A recursor is deliberately
not among them.
The Lean name is \code{KnownHead}. The plain-constant reading of the model is deliberately
\emph{not} a conjunct of the third arm: it needs the classification
Definition~\ref{def:UpstreamAsks} carries, together with the two exclusions the run supplies
at the head it is visiting.
\end{definition}

\begin{definition}[The shape condition]
\label{def:SupportedTm}
\lean{LeanToLambdaBox.SupportedTm, LeanToLambdaBox.PeanoReady, LeanToLambdaBox.peanoReadyB}
\leanok
\uses{def:PlainHead, def:KnownHead, def:CtorSaturated, def:IsLamTelescope,
def:InformativeInd, def:IndArity, def:isCasesOnName, def:Witness-SourceTable}
$\mathrm{SupportedTm}\ env\ tbl\ e\ args$ --- "$e$, read under the application spine $args$,
is inside the fragment" --- has twelve arms. Bound variables, free variables, sorts and Pi
types are unconditional (the last two are erased before they are visited). Metadata holds
\emph{only at the empty spine}, which is what makes a metadata-headed application an
exclusion, since the head-of-application function does not see through metadata. A lambda and
a \code{let} check their subterms at the empty spine. A projection carries the structure's
tabling, its informativity (restriction N18's projection half, without which the emitted
projection node is stuck on the target), the model's block data for a single-constructor
block, and the field bound. An application checks the argument on its own and reads the head
in the extended spine. A natural-number literal carries both the model's and the table's
ability to build a peano tower, pinning the two kernel constructor indices the eraser's peano
arm rebuilds. A plain constant head carries plainness, a negative \code{casesOn} test, a
negative recursor test (N21), saturation (N19) and a known head. A \code{casesOn} head
carries plainness, a positive \code{casesOn} test, the tabled inductive recovered from the
head's name prefix, its informativity, the arity bound
$\mathit{nparams} + 1 + \mathit{nindices} + 1 + |\mathrm{ctors}| \le |args|$ --- N19's
eliminator half, which the minors-length equation does \emph{not} imply at a
constructor-free inductive --- the minors slice, and a manifest lambda-telescope of the right
depth for each minor. There is deliberately \emph{no} rule for a string literal and none for
a metavariable: those shapes are outside the fragment, which is what having no rule says.
The Lean name is \code{SupportedTm}, with the peano-tower conditions \code{PeanoReady} and
\code{peanoReadyB}.
\end{definition}

\begin{remark}
The spine index is load-bearing. The checker and the relation treat metadata asymmetrically
by design: the shape function accepts a metadata node only at an empty spine accumulator, so
a metadata node under an application is excluded while one at the root is transparent. A
rendering that drops the spine index makes the metadata rule wrong.
\end{remark}

\begin{definition}[The fragment]
\label{def:Supported}
\lean{LeanToLambdaBox.Supported, LeanToLambdaBox.Reaches}
\leanok
\uses{def:SupportedTm, def:Kername, def:Witness-SourceTable}
$\mathrm{Supported}\ env\ tbl\ e$ has three fields: $e$ satisfies the shape condition at the
empty spine; so does every tabled body delta-reachable from $e$, where reachability is the
least relation containing the constants $e$ names and closed under "named by the tabled body
of a reachable name"; and no two tabled names share a \lbox{} key. It is \textbf{the coverage
predicate of the whole development}: the capstone's binder \code{hsup} and a binder of every
rung.
The Lean names are \code{Supported} and the reachability relation \code{Reaches}.
\end{definition}

\begin{remark}
The key-separation clause exists because the kername translation is \emph{not} injective
(finding F-KERNAME, refuted by Lemma~\ref{lem:toKername_not_injective}), so without it a
second tabled constant could shadow the first in the emitted environment. The Boolean that
decides the clause checks more than the clause records: it ranges over the constant column
\emph{and} the inductive column, while the recorded clause quantifies only over the constant
column. The Boolean is the stronger side, so this is not a defect --- but the fragment
predicate alone does not exclude a collision between an inductive's key and a constant's key.
\end{remark}

\begin{definition}[The safety column of the reified table]
\label{def:TableSafe}
\lean{LeanToLambdaBox.TableSafe}
\leanok
\uses{def:Witness-SourceTable, def:CtorOf}
Five clauses saying that every declaration the reified table pins is \code{safe} in the
elaboration environment --- none is \code{unsafe} or \code{partial} --- for the constant
column, for the inductive column and for that column's constructors; that no tabled constant
is an \code{\_unsafe\_rec} companion name; and that the table's two columns agree about what
a constructor is.
The Lean name is \code{TableSafe}.
\end{definition}

\begin{remark}
This is a \emph{hypothesis}, never proved: the capstone binder \code{hsafe} and a binder of
every rung, of the class whose discharge would require reading a \code{Lean.Environment},
which no Lean term denotes. The first three clauses are mechanised by
\code{lake exe reify --check}; the last two hold by construction of the reifier and no
\code{reify} verb reads them. The predicate exists because the table's adequacy pins level
parameters, types, bodies and the constructor split but \emph{not} the safety flag, while
every statement that puts a tabled name in the model reads it through the erasure
specification's declaration-adequacy field, which is stated at safe declarations only. The
\code{\_unsafe\_rec} clause excludes the name shape at which Lean's compiler lookup answers
with the \emph{other} declaration's block (finding F-UNSAFEREC). See the assumption node for
\code{hsafe} in Chapter~\ref{chap:trust}.
\end{remark}

\begin{definition}[The fix blocks the eraser installs]
\label{def:TableBlocks}
\lean{LeanToLambdaBox.TableBlocks}
\leanok
\uses{def:Erasable, def:Witness-compilerInfo, def:Witness-SourceTable}
Three clauses about the mutual blocks for which the eraser installs a fixvar map: every
member is tabled; every member's tabled body is lambda-headed, which is what makes the
emitted member body a lambda and the installed block a well-formed fixpoint node; and no
member is erasable, stated as the negation the relevance oracle's soundness contradicts.
The Lean name is \code{TableBlocks}.
\end{definition}

\begin{remark}
This is a \emph{hypothesis} --- the capstone binder \code{hblk}, of the same class as
\code{TableSafe}. The capstone does forward it to the bridge theorem, but no theorem in the
closure spends it: its only consumer in the tree is a single lemma about installed blocks
that itself has no consumer. The first two clauses are mechanised by
\code{lake exe reify --blocks} (measured: 63 blocks, all singletons, no untabled members, no
non-lambda bodies); the third is model-side and no computation reaches it. See the assumption
node for \code{hblk} in Chapter~\ref{chap:trust}.
\end{remark}

\begin{definition}[The elaborator's branch metadata]
\label{def:CasesInfoAgrees}
\lean{LeanToLambdaBox.CasesInfoAgrees, LeanToLambdaBox.CasesInfoAgrees.of_pinned}
\leanok
\uses{def:Witness-SourceTableAdequate}
Six clauses saying that the \code{Lean.CasesInfo} the shipping eraser reads for a head agrees
with the table's reified inductive: the declared name, the discriminant position, the arity,
the alternatives' range, the number of alternatives and each alternative's field count. The
companion theorem transports the clauses to the table from the kernel-side agreement, by the
table's adequacy.
The Lean names are \code{CasesInfoAgrees} and \code{CasesInfoAgrees.of\_pinned}. The
alternatives-count clause is explicitly what refutes the two early exits of the eraser's
parallel loop over alternatives.
\end{definition}

\section{The fragment checker and its soundness}
\label{sec:source:checker}

\begin{definition}[The fragment checker]
\label{def:supportedB}
\lean{LeanToLambdaBox.supportedB, LeanToLambdaBox.supportedHead, LeanToLambdaBox.supportedGo,
LeanToLambdaBox.supportedTerm, LeanToLambdaBox.constNames, LeanToLambdaBox.stepNames,
LeanToLambdaBox.reachNames, LeanToLambdaBox.saturatedB, LeanToLambdaBox.checkNames,
LeanToLambdaBox.kernameSepB, LeanToLambdaBox.informativeB, LeanToLambdaBox.resultSort,
LeanToLambdaBox.succSortB}
\leanok
\uses{def:SupportError, def:IsLamTelescope, def:CtorSaturated, def:Witness-SourceTable}
$\mathrm{supportedB}\ tbl\ \mathit{fuel}\ e$ returns \code{ok} or a support error: it first
decides key separation on the whole table, then shape-checks $e$, then unfolds the delta
closure of $e$'s constants $\mathit{fuel}$ times and, if that closure has saturated,
shape-checks every tabled body in it --- otherwise it reports exhaustion. It is
\textbf{total and fuel-indexed}: the dependency closure is cyclic through mutual blocks, so
it is not structural in $e$, and a \code{partial def} would be kernel-opaque, which is
exactly what a kernel-level \code{decide} discharge cannot afford. Exhaustion is an
\emph{error}, never a certificate. The head check decides, in order, the four plainness
tests, then either the eliminator branch (tabled inductive, informative, arity bound, one
manifest telescope per minor) or the constant branch (saturation, no recursor head, a known
column).
The Lean name is \code{supportedB}, built from the head, term and closure functions named
alongside it.
\end{definition}

\begin{remark}
The Boolean relevance test --- the declared result sort never evaluates to \code{Prop} --- is
the twin of the relation's informativity predicate, so the checker and the erasure relation
accept the same type formers; the strictly stronger syntactic successor-sort test is read
only by the first-order layer and would have rejected 9 of the 21 projection heads on the
benchmark corpus. Two anonymous \code{example}s in the module, computed by a kernel-level
\code{decide} because Lean's string-prefix test does not reduce in the elaborator, check that
the sparse-\code{casesOn} shape is reported by name at every table and that N18's projection
half is live (hand-built, because no tracked program carries such a node). Being anonymous
examples, they have no citable Lean name.
\end{remark}

\begin{lemma}[The checker's relevance test is sound]
\label{lem:informativeInd_of_tabled}
\lean{LeanToLambdaBox.informativeInd_of_tabled, LeanToLambdaBox.vResultSort_of_trExprS,
LeanToLambdaBox.ErasureSpec.contains_of_find}
\leanok
\uses{def:TableSafe, def:supportedB, def:ErasureSpec, def:Witness-SourceTableAdequate,
def:InformativeInd}
A tabled inductive type that the \emph{checker} calls informative is informative \emph{in the
model}: given an adequate, safe table over an elaboration environment that the erasure
specification connects to the model, a \code{true} verdict of the Boolean relevance test
implies the model-side informativity of the type former.
The Lean name is \code{informativeInd\_of\_tabled}.
\end{lemma}
\begin{proof}
\leanok
\uses{def:TableSafe, asm:ErasureSpec-decl_adequate}
The table pins the inductive's declared type to the elaboration environment's; safety is what
lets the erasure specification's declaration-adequacy field translate that type; and the
level translation is a homomorphism, so lean4lean's never-zero lemma carries "the result sort
never evaluates to \code{Prop}" across the translation.
\end{proof}

\begin{lemma}[A tabled inductive is the model's type former]
\label{lem:indInfo_of_tabled}
\lean{LeanToLambdaBox.indInfo_of_tabled, LeanToLambdaBox.peanoReady_of_tabled}
\leanok
\uses{def:IndInfo, def:ErasureSpec, def:Witness-SourceTableAdequate}
A tabled inductive type is the model's type former for its own block, at the parameter count
and the per-constructor field counts the table records; and the model can build a peano tower
whenever the table can.
The Lean names are \code{indInfo\_of\_tabled} and \code{peanoReady\_of\_tabled}; the second
additionally takes the table's safety.
\end{lemma}
\begin{proof}
\leanok
\uses{asm:ErasureSpec-block_adequate}
The table's adequacy pin names the block, the type's own position in it and the kernel's
indexing of its constructors; the block-adequacy field of the erasure specification reads
that as the model's block data.
\end{proof}

\begin{theorem}[Soundness of the fragment checker]
\label{thm:supportedB_sound}
\lean{LeanToLambdaBox.supportedB_sound, LeanToLambdaBox.supportedHead_sound,
LeanToLambdaBox.supportedGo_sound, LeanToLambdaBox.checkNames_sound,
LeanToLambdaBox.mem_of_reaches, LeanToLambdaBox.kernames_of_kernameSepB}
\leanok
\uses{def:Supported, def:supportedB, def:TableSafe, def:ErasureSpec,
def:Witness-SourceTableAdequate}
An \code{ok} verdict of the checker, on a table that is an adequate and safe copy of the
elaboration environment's slice, puts the subject \emph{and every delta-reachable tabled
body} inside the fragment: $\mathrm{supportedB}\ tbl\ \mathit{fuel}\ e = \code{ok}$ implies
$\mathrm{Supported}\ env\ tbl\ e$.
The Lean name is \code{supportedB\_sound}; it takes the erasure specification, the table's
adequacy and the table's safety, all three named, and is stated for every fuel.
\end{theorem}
\begin{proof}
\leanok
\uses{def:Supported, def:SupportedTm, def:supportedB, def:KnownHead, def:PlainHead,
def:CtorSaturated, def:IsLamTelescope, lem:informativeInd_of_tabled, lem:indInfo_of_tabled}
Unfold the checker. Key separation decides the third field of the fragment, by list
combinatorics on the deduplicated key list. The subject's own field is a structural induction
on \code{Lean.Expr} mirroring the checker's shape function, whose constant case is the full
case split through the head decision tree: the plainness structure is built from the four
negative tests, and then either the eliminator arm --- with
Lemma~\ref{lem:informativeInd_of_tabled} supplying informativity and the telescope decision
correspondence supplying the minors --- or the constant arm, with the known-head column built
from Lemma~\ref{lem:indInfo_of_tabled}, from the erasure specification's constructor reading,
or from "a safe declaration of the elaboration environment is a constant of the model". The
bodies field follows because a \emph{saturated} closure list holding the subject's own
constants holds everything delta-reachable from it, and every tabled body of a checked name
passed the shape check.
\end{proof}

\begin{remark}
This is the hinge from a computed verdict to the coverage hypothesis of the capstone. Every
rung of Chapter~\ref{chap:capstone} discharges its coverage binder by applying this theorem
to a \code{decide}-style computation on the committed reified table.
\end{remark}

\begin{lemma}[The fragment is closed under opening a binder]
\label{lem:SupportedTm-instantiate1}
\lean{LeanToLambdaBox.SupportedTm.instantiate1, LeanToLambdaBox.SupportedTm.instantiate1'}
\leanok
\uses{def:SupportedTm}
If a term is inside the fragment under a spine, so is the term obtained by instantiating its
outermost bound variable with a free variable --- with the spine instantiated too, because an
application's arguments are read in it.
The Lean names are \code{SupportedTm.instantiate1} and its primed form; the unprimed
statement is the primed one transported along lean4lean's modelling equation for Lean's own
instantiation function, which is why its measured axiom footprint contains that reflection
axiom.
\end{lemma}
\begin{proof}
\leanok
\uses{def:SupportedTm, asm:lean-reflection-axioms}
Induction on the derivation: every arm's side conditions are about names, tabled data or the
model and are untouched by the substitution, and the shapes commute with it. The transport to
Lean's own instantiation is a rewrite by the modelling equation.
\end{proof}

\begin{remark}
This is the form the refinement's binder cases recurse in: the eraser's monocular lambda and
let steps call the continuation on the body instantiated with a fresh free variable.
\end{remark}

\begin{lemma}[Reading the fragment back]
\label{lem:Supported-projInfo}
\lean{LeanToLambdaBox.Supported.projInfo, LeanToLambdaBox.SupportedTm.proj_inv,
LeanToLambdaBox.Supported.projInto, LeanToLambdaBox.Supported.projInd,
LeanToLambdaBox.SupportedTm.head, LeanToLambdaBox.Supported.head}
\leanok
\uses{def:SupportedTm, def:Supported, def:ProjInfo, def:IndArity}
Three read-backs the consumers need. Every projection head of a fragment term carries its
block data; in particular the informativity that the erasure relation's projection rule
demands can be recovered from a fragment verdict. And the head of a fragment term whose head
is not a constant application is itself inside the fragment at the empty spine --- which is
true precisely because of the metadata rule.
The Lean names are \code{Supported.projInfo} and the five companions.
\end{lemma}
\begin{proof}
\leanok
\uses{def:SupportedTm}
Induction on the shape derivation. At the projection arm the block-data premise yields the
model's inductive information; the other arms either have no projection head or pass the
claim to a subterm.
\end{proof}

\begin{remark}
A naming defect worth one sentence: \code{Supported.projInfo} takes a shape derivation, not a
fragment verdict, despite its namespace, so dot-notation on a fragment value does not
elaborate and its call site spells the projection out.
\end{remark}

\section{Specification environments of a run}
\label{sec:source:specenv}

\begin{definition}[A specification environment for a run state]
\label{def:SpecEnv}
\lean{LeanToLambdaBox.SpecEnv}
\leanok
\uses{def:SpecContent, def:IndCovered, def:ElimDecl, def:Kername, def:Erasure-EraseM}
$\mathrm{SpecEnv}\ env\ bo\ s\ \Gamma_{\mathrm{spec}}$ says that the \lbox{} environment
$\Gamma_{\mathrm{spec}}$ is adequate for the erasure run state $s$, in four fields: its
entries say of the source what the content predicate asks; every constant the run registered
is declared in it at its canonical kername; every inductive the run registered is
\emph{covered} --- block entry, plus, for an \emph{informative} inductive only, an eliminator
entry; and the specification bodies mention no free variable, which is what makes the
lowering pass commute with abstraction. The state occurs only through the \emph{domains} of
the two registries.
The Lean name is \code{SpecEnv}. Relevance guards the eliminator entry because the emitted
branch node of a non-informative inductive is stuck at every flag point, so relevance is a
hypothesis of both the coverage predicate and the environment relation.
\end{definition}

\begin{lemma}[Antitone in the run state]
\label{lem:SpecEnv-mono}
\lean{LeanToLambdaBox.SpecEnv.mono}
\leanok
\uses{def:SpecEnv, def:Erasure-StateLe}
An environment adequate for a state is adequate for every state below it.
The Lean name is \code{SpecEnv.mono}.
\end{lemma}
\begin{proof}
\leanok
\uses{def:SpecEnv}
The content and free-variable fields do not mention the state; the two coverage fields
compose with the two domain clauses of the state ordering.
\end{proof}

\begin{remark}
This is the property that lets sibling sub-runs be threaded under \emph{one} universally
quantified specification environment instead of merging environments. A run \emph{has} such
an environment --- \code{SpecEnv.exists} derives one from the registration invariant of
Chapter~\ref{chap:coldstart} --- so it is recovered from the run's own bookkeeping rather
than posited.
\end{remark}

\begin{lemma}[From a state's environment to a program's]
\label{lem:SpecEnv-erasesEnv}
\lean{LeanToLambdaBox.SpecEnv.erasesEnv}
\leanok
\uses{def:SpecEnv, def:ErasesEnv, def:ReachableFrom, def:Erases, def:ConstOrigin}
A specification environment of a state is a specification environment of a \emph{program} ---
\ie{} satisfies the environment relation the simulation consumes --- given three clauses the
state does not record: every kername reachable from the program is declared; every reached
compiler body has its erasure stored at \emph{every} level instantiation; and the compiler
table defines only plain constants.
The Lean name is \code{SpecEnv.erasesEnv}; the three clauses are explicit premises.
\end{lemma}
\begin{proof}
\leanok
\uses{def:SpecEnv, lem:SpecContent-erasesEnv}
Immediate from the content field and the corresponding theorem about content environments:
the clauses whose trigger is the program's reachability get it from the declaredness premise;
the other two premises are passed through, one because the content field records a single
level scope where the relation asks for every one, the other because it is not a fact about
the environment at all.
\end{proof}

\begin{definition}[The run registered everything the environment declares]
\label{def:RegSaturated}
\lean{LeanToLambdaBox.RegSaturated}
\leanok
\uses{def:ElimDecl, def:IndInfo, def:Erasure-EraseM}
Two fields: every declared non-runtime-key definition of the specification environment is
some registered constant's, and every declared block is some registered inductive's. This is
what turns the registry-\emph{scoped} clauses of a run's invariant into the unscoped clauses
of the lowering-side environment relation; without it a pruned-away definition satisfies the
scoped clause vacuously.
The Lean name is \code{RegSaturated}. It is a hypothesis at every consumer --- the capstone's
saturation binder.
\end{definition}

\begin{lemma}[Why the delta column needs parameter-free bodies]
\label{lem:defns_needs_paramFree}
\lean{LeanToLambdaBox.defns_needs_paramFree}
\leanok
\uses{def:Erases}
The delta clause's quantifier over \emph{every} level instantiation is unsatisfiable for a
body carrying its own level parameter: no erasure derivation exists for a sort at an
instantiation that sends the parameter to a level the reading scope does not know. Hence the
cold-start theorem that discharges the clause does so at level-parameter-free bodies and no
further.
The Lean name is \code{defns\_needs\_paramFree}; the Lean statement is the negation of the
quantified clause at the sort built from a level parameter, for an arbitrary target term.
\end{lemma}
\begin{proof}
\leanok
\uses{def:Erases, asm:lean-reflection-axioms}
Instantiate at a substitution sending the parameter to a level not in the reading scope; the
instantiation is computed by the modelling equation for Lean's level substitution, which is
where the measured reflection axioms enter. The only erasure arm applicable to a sort is the
box arm, whose translation witness needs the level in scope; inverting it gives a
level-translation equation that is \code{none}, a contradiction.
\end{proof}

\begin{remark}
A refutation that pins the shape of a hypothesis rather than proving a positive fact is
exactly the kind of self-audit a blueprint should record.
\end{remark}

\begin{proposition}[The chain is not vacuous]
\label{prop:lowerEnv_of_cold_run}
\lean{LeanToLambdaBox.lowerEnv_of_cold_run, LeanToLambdaBox.fvarFree_idEnv,
LeanToLambdaBox.demoEnv_specEnv, LeanToLambdaBox.fvarFree_demoEnv}
\leanok
\uses{def:LowerEnv, def:SpecContent, def:SpecEnv, def:RegSaturated}
A run that starts at the empty state and registers one lambda-bodied definition satisfies the
registration invariant and saturates its own specification environment, so the lowering-side
environment relation comes out with \emph{every} clause derived. A four-entry demo fixture is
likewise a specification environment for the initial state.
The Lean name is \code{lowerEnv\_of\_cold\_run}; its one premise is that the one-entry
environment satisfies the content predicate.
\end{proposition}
\begin{proof}
\leanok
\uses{thm:RegInvShape-prime-lowerEnv, def:RegSaturated}
Build the registration invariant by one registration step from the empty-state introduction;
then discharge the two saturation clauses on the one-entry environment by inspecting the
lookup, and apply the cold-start theorem that converts an invariant plus saturation into the
lowering-side relation.
\end{proof}

\begin{proposition}[Eliminator coverage, at the Nat fixture]
\label{prop:natEnv_elimCovered}
\lean{LeanToLambdaBox.natEnv_elimCovered, LeanToLambdaBox.natSpecEnv,
LeanToLambdaBox.natSpecMib, LeanToLambdaBox.natSpecEnv_block,
LeanToLambdaBox.natSpecEnv_elim}
\leanok
\uses{def:ElimDecl, def:IndInfo, def:mkElimBody, def:IndCovered}
For the hand-built \code{Nat} fixture, the two-entry \lbox{} environment holding the block and
the eliminator's canonical body supplies exactly what eliminator coverage returns: the
eliminator declaration at the \code{casesOn} constant's kername, the model's block data
beside it, and the constructor count.
The Lean name is \code{natEnv\_elimCovered}; the Lean statement is the three-fold conjunction
spelled out at the fixture's own identifiers.
\end{proposition}
\begin{proof}
\leanok
\uses{prop:NatWitness-seval_iota_fires}
The two lookups are the head and the one-key-past case of environment lookup; the model side
is the fixture's own block data.
\end{proof}

\begin{remark}
This is the non-vacuity witness for the eliminator clause of the coverage predicate.
\end{remark}

\section{What lean4lean still owes}
\label{sec:source:upstream}

\begin{definition}[A block declared below the environment]
\label{def:IndBlockBelow}
\lean{LeanToLambdaBox.IndBlockBelow, LeanToLambdaBox.Peel}
\leanok
\uses{def:lean4lean-grounding}
A mutual inductive declaration is \emph{below} an environment when some well-formed
declaration list containing it produces an environment at or under it. Every classification
predicate of this development exhibits its block in this form, which is why the uniqueness ask
must be stated in it. Separately, \emph{peeling} a spine's typing argument by argument says
the head's type is definitionally a Pi whose domain types the first argument and whose
instantiated codomain peels the rest.
The Lean names are \code{IndBlockBelow} and \code{Peel}. The first-order layer's own block
predicate is the special case where the declaration list's environment is the environment
itself.
\end{definition}

\begin{definition}[The load-bearing upstream asks, as one premise]
\label{def:UpstreamAsks}
\lean{LeanToLambdaBox.UpstreamAsks}
\leanok
\uses{def:IndBlockBelow, def:CtorOf, def:IndInfo, def:ConstOrigin, def:CasesOnShape, def:lean4lean-grounding, asm:UpstreamAsks-constsOrigin, asm:UpstreamAsks-constArityInv, asm:UpstreamAsks-mkAppsInv, asm:UpstreamAsks-indSpineInj}
One \code{Prop}-structure with four fields, packaging the four facts about lean4lean's kernel
theory that this development needs and cannot prove here (\code{doc/upstream-asks.md} items
2, 6, 9 and 10).
\begin{itemize}
\item The constants-keyed origin field, in eight conjuncts: a constructor is not an inductive
      type name; the constructor reading determines its type and index; the block reading
      determines its identifier, parameter count and field counts; the eliminator reading
      determines its segmentation; every declared constant is a constructor, a type name or a
      plain constant; the plain-constant reading excludes the other two; a block \emph{below}
      the environment declaring a type former is matched by a block of the environment's own
      declaration list; and two blocks below the environment declaring the same type former
      are the same block.
\item The arity-inversion field: an application headed by an inductive type former is
      definitionally equal to neither a sort nor a Pi.
\item The spine-typing-inversion field, stated with the strong-ordering premise explicit, so
      that the ask is \code{sorryAx}-free exactly as the corresponding application-inversion
      lemma is.
\item The spine-injectivity field: two definitionally equal spines headed by
      \emph{inductively declared} type formers have the same head. The declaredness premises
      are load-bearing, since without them an ordinary definition refutes the statement.
\end{itemize}
The Lean name is \code{UpstreamAsks}.
\end{definition}

\begin{remark}
\textbf{Proof status: hypothesis. Nothing under this structure is proved in this repository.}
The four fields are the four \code{UpstreamAsks} assumption nodes of
Chapter~\ref{chap:trust} (Assumptions~\ref{asm:UpstreamAsks-constsOrigin},
\ref{asm:UpstreamAsks-constArityInv}, \ref{asm:UpstreamAsks-mkAppsInv} and
\ref{asm:UpstreamAsks-indSpineInj}). The structure is taken by the
erasure-correctness theorem, by the bridge, by the capstone and by every green rung, and a
single pin bump would discharge all four with no change to any consumer's statement shape.
The repository's trust document carries the row and the \emph{measured} ledger cannot, because
\code{\#print axioms} reports the footprint of a proved term and never a hypothesis. Two asks
are deliberately \emph{not} fields (the defining-equation ownership lemma and the seven
kernel-generic declarations), having no consumer here; and the first conjunct of the first
field is proved in this repository by Lemma~\ref{lem:CtorOf-not_indInfo}, so in a dependency
graph a proved fact sits inside a hypothesis node.
\end{remark}

\section*{Deviations from the paper and caveats}

\textbf{There is no source semantics to cite, so one is defined.} $\SEval$ is not a
free-standing invention: it is a presentation of a sub-relation of the kernel's own
convertibility, with the evaluation order made explicit and every non-value arm paying a
definitional-equality obligation. Where [S] can quote PCUIC reduction, this development must
build and justify its analogue, which is what Section~\ref{sec:source:values} does.

\textbf{Delta is not kernel delta} (tag N8). The eraser reads the compiler's bodies --- after
\code{csimp} rewriting and \code{@[extern]} or \code{@[implemented\_by]} substitution --- so
the delta arm unfolds a table independent of the kernel environment and pays with a
definitional-equality obligation at the applied instance. The corresponding typing hypothesis
(Definition~\ref{def:CompilerBodies}) is a capstone binder discharged at the first rung only.

\textbf{The coverage predicate has no counterpart in [S].} MetaCoq's erasure is total on
well-typed PCUIC; a coverage predicate exists here because the subject is a shipping, partly
panicking eraser read against a \emph{compiler} body table. In [L]'s terms it is the explicit
statement of which part of the extraction's input domain the correctness argument covers. Of
the five benchmark programs, only Arith is in the fragment. The recursor-head exclusion
prevents \emph{vacuous} coverage, since a recursor spine has no source evaluation at all;
the \code{@[implemented\_by]} hole is closed not by the fragment but by the configuration
hypothesis \code{ConfigPinned}, which is \emph{not} the default \code{\#erase} configuration
--- it differs from it on three of its five fields.

\textbf{Specification environments are \code{erases\_global} derived from a run}, not posited:
the eraser registers constants and blocks as it discovers them, so the environment relation is
recovered from the run's own bookkeeping.

\textbf{Fixture caveats.} Two of the three witness environments are proved well-formed only
\emph{before} the hand-added conversion equations they need, and the delta fixture's
environment is never shown well-formed at all; the right reading of all three is "the arms'
side conditions are jointly satisfiable".
